% Copyright (c) 2023-2026 Oriol Cayon, Delft University of Technology
% SPDX-License-Identifier: Apache-2.0
\documentclass[11pt,a4paper]{article}

\usepackage[margin=2.4cm]{geometry}
\usepackage{graphicx}
\usepackage{amsmath,amssymb}
\usepackage{booktabs}
\usepackage{longtable}
\usepackage{xcolor}
\usepackage{caption}
\usepackage{float}
\usepackage{enumitem}
\usepackage{fancyhdr}
\usepackage[hidelinks]{hyperref}

\graphicspath{{../../results/structural/demo/}{../../results/structural/validation/}}

\definecolor{accent}{HTML}{D55E00}
\definecolor{muted}{HTML}{555555}

\captionsetup{font=small,labelfont=bf}
\setlist{itemsep=2pt,parsep=0pt,topsep=3pt}

\pagestyle{fancy}
\fancyhf{}
\fancyhead[L]{\small\textsc{Billow}}
\fancyhead[R]{\small\thepage}
\renewcommand{\headrulewidth}{0.3pt}

\newcommand{\code}[1]{\texttt{\small #1}}

\title{%
  \vspace{-1.5cm}
  {\Huge\textbf{Billow}}\\[6pt]
  {\large A minimum-energy structural model for soft kites}\\[10pt]
  {\normalsize Timoshenko beams, inflatable tubes, wrinkling membranes,\\
   cables and pulleys --- one potential energy, one NLP}
}
\author{Oriol Cayon \\ \small Delft University of Technology}
\date{\today}

\begin{document}
\maketitle
\thispagestyle{fancy}

\begin{abstract}
\noindent
Billow is a standalone structural library for soft-wing airborne wind energy
systems. It assembles cables, frictionless pulleys, geometrically exact
Timoshenko beams, inflatable tubes and wrinkling membrane fabric into a single
total potential energy, and finds static equilibrium by minimising it with
IPOPT. Stationarity of the energy \emph{is} the equilibrium equation, so a
converged solve satisfies the nodal force balance to solver tolerance rather
than to whatever residual a relaxation scheme happens to leave behind.

This document specifies the formulation, the element library, the validation
against benchmarks external to the code, a set of demonstration cases, and a
measured comparison against \code{kite\_fem}, the existing FEM path. Headline
results: exact-to-$10^{-4}$ agreement with the Euler elastica, second-order
convergence to the exact roll-up circle, agreement with published values for
the Bathe--Bolourchi $45^\circ$ bend, reproduction of the ASKITE inflatable-tube
fits to $0.06\%$ in bending, and a full-kite structural solve in $5.9$\,s cold
and $0.58$\,s warm at 1893 degrees of freedom.

Billow currently lives as \code{awetrim.structural} and imports only NumPy and
CasADi. It is intended to become its own repository.
\end{abstract}

\tableofcontents
\newpage

%======================================================================
\section{Scope and design intent}

\subsection{What Billow is}

A static structural solver for the kinds of structures soft kites are made of:
thin lines that cannot push, ropes that run over sheaves, inflated tubes that
bend until they collapse, and fabric that carries tension but wrinkles instead
of carrying compression. Given a geometry, material properties and a set of
nodal loads, it returns the deformed shape and the internal force distribution.

\subsection{What Billow is not}

It contains \textbf{no aerodynamics}. There are no polars, no airfoils, no
lift model. Loads arrive as nodal forces from outside. Everything in this
document that involves pressure uses a prescribed uniform differential
pressure, not an aerodynamic solution. The same is true of \code{kite\_fem},
which is also purely structural.

It is also not a dynamics code: there are no inertial terms and no time
integration. Equilibrium is quasi-static.

\subsection{Isolation}

Billow depends only on NumPy and CasADi. It does not import the particle-system
solver (PSS), the vortex-step method (VSM), the aerostructural coupling loop or
any kite YAML schema, and nothing in those modules imports it. This is
deliberate and is the first invariant to preserve: it allows the element physics
to be validated against closed-form solutions before any of it touches a coupled
run. A coupling adapter, when written, belongs on the aerostructural side, not
inside Billow.

\subsection{Why a new model}

The existing minimum-energy solver in the AWETrim tree
(\code{aerostructural/pss/structural\_nlp.py}) already solves the current bridle
well. It does not extend, for two independent reasons.

\begin{description}[style=unboxed,leftmargin=1.2em]
\item[Physics.] Its element vocabulary is axial springs. Battens and
leading-edge tubes need bending and torsion; the canopy needs a membrane that
cannot carry compression.
\item[Scale.] It builds the objective with a Python loop over elements in
CasADi \code{MX}, so the expression graph grows with the element count. At tens
of springs this is free. At tens of thousands of fabric triangles, graph
construction --- and especially \code{gradient} and \code{hessian} over it ---
becomes the bottleneck long before the numerical solver does.
\end{description}

Billow keeps the formulation (still an NLP, still IPOPT, still minimum potential
energy) and replaces the assembly.

%======================================================================
\section{Formulation}

\subsection{Total potential energy}

The unknown configuration $X$ minimises
\begin{equation}
  \Pi(X) \;=\; \sum_{\text{sets}} U_{\text{set}}(X)
  \;-\; \mathbf{f}_{\text{ext}}\!\cdot\!\mathbf{x}
  \;-\; \mathbf{m}_{\text{ext}}\!\cdot\!\boldsymbol{\theta}(\boldsymbol{\psi})
  \;+\; \tfrac{1}{2}k_a\lVert \mathbf{x}-\mathbf{x}_{\text{seed}}\rVert^2 ,
  \label{eq:pi}
\end{equation}
subject only to bounds on pinned degrees of freedom. The last term is a weak
anchor discussed in \S\ref{sec:anchor}.

Because $\partial\Pi/\partial X = 0$ is the equilibrium equation, a converged
solve satisfies $\mathbf{f}_{\text{int}} + \mathbf{f}_{\text{ext}} = 0$ at every
free node to solver tolerance. There is no separate residual criterion to tune.

\subsection{Degrees of freedom}

\begin{equation}
  X = \big[\; \underbrace{\mathbf{x}_0 \ldots \mathbf{x}_{n-1}}_{\text{3 per node}}
  \;\big|\;
  \underbrace{\boldsymbol{\psi}_0 \ldots \boldsymbol{\psi}_{r-1}}_{\text{3 per rotational node}} \;\big]
\end{equation}

Three translations for \emph{every} node, and three incremental rotations only
for the $r$ nodes that some beam element needs a material frame for. Cable,
pulley and membrane nodes stay at three degrees of freedom. A canopy of 10\,000
fabric nodes therefore costs 30\,000 unknowns, not 60\,000, regardless of what
beams are attached elsewhere. For contrast, \code{kite\_fem} allocates six
degrees of freedom per node everywhere and masks the unused ones.

\subsection{Finite rotations: the Cayley parameterisation}

Rotations use the \textbf{Rodrigues (Cayley) vector}
$\boldsymbol{\psi} = 2\tan(\phi/2)\,\mathbf{n}$ rather than the exponential
map's rotation vector $\phi\,\mathbf{n}$. With $S = \operatorname{skew}(\boldsymbol{\psi})$,
\begin{align}
  R(\boldsymbol{\psi}) &= I + \frac{4}{4+\boldsymbol{\psi}\cdot\boldsymbol{\psi}}
    \left(S + \tfrac{1}{2}S^2\right), \label{eq:cayley}\\
  \boldsymbol{\psi}(R) &= \frac{2\,\operatorname{axial}(R-R^{\mathsf T})}{1+\operatorname{tr}R},
    \label{eq:cayleyinv}\\
  \boldsymbol{\psi}_{1/2} &= \frac{\boldsymbol{\psi}}
    {1+\sqrt{1+\tfrac{1}{4}\boldsymbol{\psi}\cdot\boldsymbol{\psi}}},
    \qquad R(\boldsymbol{\psi}_{1/2})^2 = R(\boldsymbol{\psi}).
    \label{eq:cayleyhalf}
\end{align}

The motivation is entirely numerical. Every map above is a \emph{rational}
function of its argument. There is no $\sin\phi/\phi$ removable singularity, so
no \code{if\_else} branch is needed at $\phi=0$, so the Hessian that IPOPT
differentiates is continuous everywhere in $|\phi| < \pi$. The half-angle
identity \eqref{eq:cayleyhalf} is worth noting: the naive form
$\tan(\phi/4) = (\sqrt{1+t^2}-1)/t$ has a removable $0/0$ at $t=0$;
rationalising it removes the singularity exactly.

\subsection{Incremental rotations and frame updates}

$\boldsymbol{\psi}_j$ is measured against a stored reference frame, so the
configuration is $R_j = R(\boldsymbol{\psi}_j)\,R_{\text{ref},j}$ and every
solve starts from $\boldsymbol{\psi} = 0$. After each solve the increment is
absorbed, $R_{\text{ref}} \leftarrow R(\boldsymbol{\psi}^\star) R_{\text{ref}}$,
and re-projected onto $SO(3)$ so repeated absorption cannot drift.

The Cayley chart is singular at $\phi = \pi$. Incremental rotations mean a node
can turn arbitrarily far \emph{across a load ramp} while each individual solve
stays well inside the chart. This is why applied moments require load stepping
(\S\ref{sec:moments}); nodal forces do not.

\subsection{Applied moments are not conjugate to \texorpdfstring{$\boldsymbol{\psi}$}{psi}}
\label{sec:moments}

A moment does work against the \emph{exponential-map} rotation vector, not
against the Rodrigues vector the degrees of freedom carry. The external work
term therefore converts,
\begin{equation}
  \boldsymbol{\theta}(\boldsymbol{\psi}) = \boldsymbol{\psi}\,
  \frac{2\arctan\!\left(\lVert\boldsymbol{\psi}\rVert/2\right)}{\lVert\boldsymbol{\psi}\rVert},
\end{equation}
which tends to $\boldsymbol{\psi}$ as $\lVert\boldsymbol{\psi}\rVert \to 0$.

This is not cosmetic. Using $-\mathbf{m}\cdot\boldsymbol{\psi}$ directly is
correct only to first order, and the roll-up benchmark
(\S\ref{sec:rollup}) exposed it as a \emph{load-path dependence}: the same
static problem gave $8.1\times10^{-3}L$ with four load steps and
$2.1\times10^{-1}L$ with eight. Statics must be path independent. With the
conversion in place the answer is identical to five digits across 4, 8 and 16
steps, and this is guarded by a regression test.

Two limitations remain, both fundamental rather than implementation artefacts:

\begin{itemize}
\item Applied moments need load stepping, because the chart is singular at
$\phi = \pi$. Four steps suffice for a full $2\pi$ roll-up.
\item A constant (``dead'') moment in three dimensions is \textbf{genuinely
non-conservative} and therefore has no potential at all. Energy minimisation can
represent only the fixed-axis case.
\end{itemize}

Nodal \emph{forces} --- which is what an aerodynamic coupling applies --- carry
none of these caveats and are exact in a single solve.

\subsection{Assembly: one mapped kernel per element type}
\label{sec:assembly}

This is the single design decision that makes the model scale. Each element
kernel is compiled once to a small CasADi \code{SX} function over \emph{one}
element and evaluated across the whole set with \code{Function.map} on a
gathered index matrix:
\[
  Q = \operatorname{reshape}\big(X[\,\mathrm{idx}\,],\; n_{\text{dof/el}},\; n_{\text{el}}\big),
  \qquad
  U_{\text{set}} = \textstyle\sum_e \mathcal{U}(Q_{:,e}, F_{:,e}, P_{:,e}).
\]
The objective graph therefore holds one node per element \emph{type}, not per
element. Graph construction, gradient and Hessian generation cost
$O(\text{number of element types})$, not $O(\text{number of elements})$.

The Hessian retains the sparsity of the element connectivity graph, so IPOPT's
sparse linear algebra sees a mesh-stencil matrix and scales accordingly. The
measured consequence is Table~\ref{tab:scaling}.

\begin{table}[H]
\centering
\caption{Flat clamped canopy, uniform out-of-plane load. Iteration count is
essentially mesh independent; cost sits in the sparse linear algebra.
Source: \code{run\_demo\_cases.py --case scaling}.}
\label{tab:scaling}
\begin{tabular}{rrrrr}
\toprule
Triangles & DOF & Build (s) & Solve (s) & IPOPT iterations \\
\midrule
128  & 243  & 0.07 & 0.05 & 6 \\
1800 & 2883 & 0.67 & 1.25 & 9 \\
5832 & 9075 & 3.0  & 6.1  & 9 \\
\bottomrule
\end{tabular}
\end{table}

\subsection{Everything numeric is a parameter}

Reference frames, every element parameter column (rest lengths, stiffnesses,
section properties), the external loads and the anchor are NLP
\emph{parameters}, not baked constants. One solver build therefore serves an
entire sweep: tape actuation, a stiffness ramp or a load sweep change parameter
values, never topology.

\subsection{The anchor term}
\label{sec:anchor}

The final term of \eqref{eq:pi} is a weak spring to the seed positions, present
only to pin a node that no taut element reaches and whose equilibrium position
would otherwise be indeterminate. It biases the reported force balance by at
most $k_a\lVert\mathbf{x}-\mathbf{x}_{\text{seed}}\rVert$; the default
$k_a = 10^{-6}$\,N/m puts that far below any coupling tolerance.

\subsection{Convergence criterion}

\code{converged} is the \emph{physical} verdict: the solve is accepted when the
free nodes are in force balance to
\[
  \max\big(\varepsilon_{\text{abs}},\; \varepsilon_{\text{rel}}\max_i\lVert \mathbf{f}_{\text{ext},i}\rVert\big),
  \qquad \varepsilon_{\text{abs}} = 10^{-6}\,\text{N},\;\; \varepsilon_{\text{rel}} = 10^{-6},
\]
whether or not IPOPT liked its own termination test. The relative term matters:
an absolute newton means very different things on a 1\,N bridle line and a
600\,N benchmark beam. IPOPT's own verdict is preserved alongside as
\code{ipopt\_success} and \code{status} --- on stiff, nearly quadratic
structures it routinely reports \code{Search\_Direction\_Becomes\_Too\_Small} at
an answer already exact to machine precision.

%======================================================================
\section{Element library}

Every kernel returns the strain energy of one element. Kernels are pure: no
loops over elements, no global indexing, no closing over per-element data.
Everything element-specific arrives through the parameter row.

\subsection{Cable}

A linear axial spring with stretch $s = \ell - \ell_0$:
\begin{equation}
  U = \tfrac{1}{2}k\,\langle s\rangle^2,
  \qquad \langle s\rangle = \max(0, s)\ \text{when tension-only}.
\end{equation}
The tension-only cut is $C^1$ but not $C^2$: the stiffness jumps across the
taut/slack boundary. An optional \code{slack\_smoothing} width rounds the corner
to recover a continuous Hessian at the cost of a small spurious compressive
stiffness.

\subsection{Pulley}

A rope running over a frictionless sheave at the middle node, with a single
shared stretch across both arms:
\begin{equation}
  U = \tfrac{1}{2}k\,\big\langle \lVert\mathbf{x}_2-\mathbf{x}_1\rVert
    + \lVert\mathbf{x}_3-\mathbf{x}_2\rVert - \ell_0 \big\rangle^2 .
\end{equation}
One shared stretch is exactly the statement that rope tension is equal either
side of the sheave, and it lets the pulley node find its own position along the
rope.

\paragraph{Convention warning.} \code{rest\_length} here is the \textbf{whole
rope}, both arms. This is not universal: the PSS geometry reader splits $\ell_0$
across the two arms, while \code{kite\_fem} stores the total on each arm. Any
geometry adapter must state which convention it is reading.

\subsection{Geometrically exact Timoshenko beam}

Two-node Simo--Reissner beam with one-point (mid-element) integration. Each node
carries a position and an orthonormal material frame $R = [\mathbf{d}_1\,
\mathbf{d}_2\, \mathbf{d}_3]$: twelve degrees of freedom per element.

With $\boldsymbol{\psi}$ the relative rotation between the two node frames from
\eqref{eq:cayleyinv}, and $R_m$ the exact half-way frame from
\eqref{eq:cayleyhalf}:
\begin{align}
  \boldsymbol{\Gamma} &= R_m^{\mathsf T}(\mathbf{x}_2-\mathbf{x}_1)/L_0
    && \text{axial stretch and two transverse shears},\\
  \boldsymbol{\Omega} &= R_m^{\mathsf T}\boldsymbol{\psi}/L_0
    && \text{torsion and two bending curvatures},\\
  U &= \tfrac{L_0}{2}\Big[
     \Delta\boldsymbol{\Gamma}^{\mathsf T}\operatorname{diag}(EA, GA_2, GA_3)\Delta\boldsymbol{\Gamma}
   + \Delta\boldsymbol{\Omega}^{\mathsf T}\operatorname{diag}(GJ, EI_2, EI_3)\Delta\boldsymbol{\Omega}
    \Big],
\end{align}
with $\Delta(\cdot) = (\cdot) - (\cdot)_0$ measured against the reference
configuration, so an initially curved or pre-twisted member is stress-free as
built.

Three properties make this element usable inside an NLP objective:

\begin{description}[style=unboxed,leftmargin=1.2em]
\item[Objective.] Under a superposed rigid rotation $Q$, both $\boldsymbol{\psi}$
and $\mathbf{x}_2-\mathbf{x}_1$ map to $Q(\cdot)$ while $R_m \to QR_m$, so
$\boldsymbol{\Gamma}$ and $\boldsymbol{\Omega}$ are invariant and rigid-body
motion costs exactly zero energy. A slack member that swings freely is not
spuriously stiffened.
\item[Shear-locking free.] The single midpoint sample is the standard reduced
integration cure. Kite battens and leading-edge tubes are slender, so this is
not optional; \S\ref{sec:locking} verifies it at slenderness 800.
\item[Branch-free.] Using the Cayley vector as the curvature measure makes the
element strain $\boldsymbol{\Omega}_{\text{exact}} + O(\phi^3)$ per element ---
the same order as the one-point quadrature error, and vanishing under mesh
refinement. This is a deliberate trade: the exponential map's $\sin\phi/\phi$
would need a branch at $\phi=0$.
\end{description}

\subsection{Inflatable tube}
\label{sec:inflatable}

An inflated leading-edge tube is not linear-elastic. Its bending stiffness falls
away as the compressed side wrinkles, until the section carries a limiting
moment and then collapses. \code{kite\_fem} captures this with empirical fits
(coefficients $C_1 \ldots C_{19}$, from the ASKITE work) applied as \emph{secant}
stiffnesses: $EI$ and $GJ$ are recomputed each iteration from the current
deflection and twist.

That is consistent inside a Newton residual, but it is \textbf{not a potential}.
Substituting a state-dependent $EI(\kappa)$ into $\tfrac{1}{2}EI\kappa^2$
silently drops the $\mathrm{d}EI/\mathrm{d}\kappa$ terms, so the gradient would
no longer be the internal moment and a minimum-energy solve would converge to
the wrong equilibrium. The fits must be \emph{integrated}, not substituted.

\paragraph{From tip-load fit to moment--curvature law.}
The bending fit is published as tip load against normalised tip deflection
$v = \delta/L$ of a \textbf{one-metre} cantilever:
$P(v) = D\big(1-e^{-(N/D)v}\big)$. For a linear cantilever
$\delta = PL^3/3EI$ and $\kappa_{\text{root}} = PL/EI$, so $\kappa = 3v/L$, and
the root moment is $M = PL$. At the calibration length $L = 1$\,m this converts
the fit into an intrinsic constitutive law:
\begin{align}
  M(\kappa) &= M_{\max}\left(1 - e^{-EI_0\kappa/M_{\max}}\right),
   \qquad EI_0 = N/3, \quad M_{\max} = D, \label{eq:tubemoment}\\
  W_b(\kappa) &= M_{\max}\left[\,|\kappa| - k_0\left(1 - e^{-|\kappa|/k_0}\right)\right],
   \qquad k_0 = M_{\max}/EI_0 . \label{eq:tubeenergy}
\end{align}
$W_b$ is quadratic ($\tfrac{1}{2}EI_0\kappa^2$) near zero, so it is smooth there
despite the $|\kappa|$, and saturates at $M_{\max}$.

Recasting this way also removes a length inconsistency: \code{kite\_fem} infers
$EI = P/(3v)$, which equals the true $EI$ only for a one-metre element, because
$v$ is already normalised by element length. A moment--curvature law is
length-independent by construction.

\paragraph{Torsion} is already intrinsic (twist per unit length) and integrates
directly:
\begin{equation}
  T(\omega) = c_1\arctan(c_2\omega),
  \qquad
  W_t(\omega) = c_1\!\left[\omega\arctan(c_2\omega)
    - \frac{\ln\!\left(1+c_2^2\omega^2\right)}{2c_2}\right],
\end{equation}
with $GJ_0 = c_1c_2$ and a limiting torque $c_1\pi/2$.

Biaxial bending uses the resultant curvature $\sqrt{\Omega_2^2+\Omega_3^2}$, the
correct isotropic extension for an axisymmetric tube. Axial and shear stiffness
are \emph{not} covered by the fits and stay linear; supply them from the tube
geometry.

\begin{table}[H]
\centering
\caption{Derived tube properties from the ASKITE fits.
Higher pressure stiffens and strengthens; the thinner tube collapses
\emph{before} its moment saturates ($\kappa_{\text{collapse}} < k_0$).}
\label{tab:tube}
\begin{tabular}{llrrrrrr}
\toprule
$d$ (m) & $p$ (bar) & $EI_0$ (N\,m$^2$) & $M_{\max}$ (N\,m) & $k_0$ (m$^{-1}$)
 & $\kappa_{\text{collapse}}$ (m$^{-1}$) & $GJ_0$ (N\,m$^2$) & $T_{\max}$ (N\,m)\\
\midrule
0.16 & 0.3 & 288.2 & 58.10 & 0.2016 & 0.2751 & 91.3 & 124.9 \\
0.16 & 0.5 & 362.0 & 98.78 & 0.2728 & 0.2973 & 86.5 & 174.7 \\
0.12 & 0.3 & 126.8 & 34.16 & 0.2694 & 0.2217 & 33.2 & 117.1 \\
\bottomrule
\end{tabular}
\end{table}

\paragraph{Collapse is reported, not enforced.}
A genuinely dropping post-collapse moment would make the energy \emph{decrease}
with curvature --- an unbounded mechanism that minimisation would simply run
away from. The law therefore stops at saturation, and
\code{inflatable\_beam\_state} returns a \code{utilisation} ratio
$\kappa/\kappa_{\text{collapse}}$ and a \code{collapsed} flag, so leaving the
calibrated range is visible rather than silent.

\subsection{Wrinkling membrane}
\label{sec:membrane}

Constant-strain triangle with a Saint-Venant--Kirchhoff plane-stress law, in the
finite-strain (Green--Lagrange) measure so a canopy panel can rotate and billow
arbitrarily without spurious stiffening. Translation-only: a membrane carries no
bending.

The reference triangle is flattened once into its own in-plane basis and stored
as the inverse edge matrix $D_0^{-1}$, so the kernel only ever sees a $2\times2$
map:
\begin{equation}
  F = [\,\mathbf{a}\;\mathbf{b}\,]\,D_0^{-1}, \quad
  C = F^{\mathsf T}F, \quad
  E = \tfrac{1}{2}(C-I),
\end{equation}
\begin{equation}
  W_{\text{taut}} = \frac{E_{\text{mod}}}{2(1-\nu^2)}
   \left[(\operatorname{tr}E)^2 - 2(1-\nu)\left(E_{11}E_{22}-E_{12}^2\right)\right].
\end{equation}

\paragraph{Wrinkling.} A membrane cannot carry compression. Left alone, the SVK
law lets the canopy mesh buckle into whatever compressive mode the
discretisation offers, and the result depends on mesh density rather than on
physics. The relaxed (tension-field) energy in principal Green strains
$e_1 \ge e_2$ is the classical Pipkin form:
\begin{equation}
  W_{\text{relaxed}} =
  \begin{cases}
    W_{\text{taut}}, & e_2 + \nu e_1 \ge 0 \quad\text{(taut)},\\[2pt]
    \tfrac{1}{2}E_{\text{mod}}\,e_1^2, & e_1 > 0 \ \text{otherwise}\quad\text{(wrinkled)},\\[2pt]
    0, & e_1 \le 0 \quad\text{(slack)}.
  \end{cases}
\end{equation}

This is exactly what a minimum-energy formulation wants: the relaxed functional
is the quasiconvex envelope, so the minimiser lands on the wrinkled state
directly instead of chasing the near-singular tangent a residual-form Newton
solve has to fight through. \textbf{Solving fabric by energy minimisation is not
a workaround --- it is the better-posed formulation.}

\paragraph{Slack fabric needs a residual stiffness.} A fully slack region stores
\emph{exactly} zero energy, so its Hessian block is exactly zero and the Newton
step is undefined; IPOPT reports \code{Error\_In\_Step\_Computation} on a panel
clamped into a frame smaller than itself. Blending a small fraction of the
unrelaxed law back in,
\begin{equation}
  W = (1-r)\,W_{\text{relaxed}} + r\,W_{\text{taut}}, \qquad r = 10^{-4},
\end{equation}
leaves the taut region untouched (there the two coincide) while giving slack
fabric a definite tangent. It is also not a fiction: real fabric has some
compressive and bending resistance. Do not set $r=0$ on a canopy that can go
slack.

%======================================================================
\section{Validation}

All references in this section come from \emph{outside} the code: an exact
boundary-value solution, an exact circle, published tabulated values, or the
independent fits themselves. Nothing here can drift with the implementation.

Where \code{kite\_fem} appears it is driven with matched linear section
properties and its inflatable constitutive update disabled, so the two share one
material model. Section properties are deliberately slender: the classical
references are inextensible and shear-rigid, so axial and shear flexibility are
kept below $\sim0.1\%$ and cannot be mistaken for discretisation error.

\subsection{Euler elastica}

A tip-loaded cantilever at load parameter $\alpha = PL^2/EI$ up to 10. The
reference is the exact elastica, obtained by solving
$\theta'' = -\alpha\cos\theta$, $\theta(0)=0$, $\theta'(1)=0$ as a
boundary-value problem to $10^{-10}$ with continuation in $\alpha$ --- no beam
element is involved.

\begin{figure}[H]
\centering
\includegraphics[width=\textwidth]{elastica.png}
\caption{Tip locus and error against the exact elastica, 40 elements.
Billow tracks the exact curve across the whole range; \code{kite\_fem}
departs progressively and its tip deflection saturates beyond $\alpha \approx 2$.}
\label{fig:elastica}
\end{figure}

\begin{table}[H]
\centering
\caption{Tip position error normalised by length, 40 elements.}
\begin{tabular}{lrrrr}
\toprule
$\alpha$ & 0.5 & 2 & 6 & 10 \\
\midrule
Billow           & $9.0\times10^{-5}$ & $2.4\times10^{-4}$ & $6.5\times10^{-4}$ & $1.1\times10^{-3}$\\
\code{kite\_fem} & $1.6\times10^{-2}$ & $1.8\times10^{-1}$ & $4.2\times10^{-1}$ & $4.9\times10^{-1}$\\
\bottomrule
\end{tabular}
\end{table}

\paragraph{Mesh refinement separates discretisation from formulation.}
At $\alpha = 3$, refining from 10 to 160 elements:

\begin{table}[H]
\centering
\begin{tabular}{lrrrrr}
\toprule
Elements & 10 & 20 & 40 & 80 & 160\\
\midrule
Billow           & $3.8\times10^{-3}$ & $7.0\times10^{-4}$ & $3.1\times10^{-4}$
                 & $4.6\times10^{-4}$ & $5.0\times10^{-4}$\\
\code{kite\_fem} & $2.908\times10^{-1}$ & $2.914\times10^{-1}$ & $2.916\times10^{-1}$
                 & $2.917\times10^{-1}$ & $2.917\times10^{-1}$\\
\bottomrule
\end{tabular}
\end{table}

The \code{kite\_fem} error is \textbf{mesh independent}, which places it in the
formulation rather than the discretisation. Its beam matches linear theory to
$0.02\%$ below about one degree of tip rotation, so the property matching used
here is correct; the error is purely rotation-driven. Envelope: below $1\%$ up
to roughly $10^\circ$, about $15\%$ at $\alpha=1$, saturating beyond $\alpha=2$.
It reports \code{converged=True} at those answers.

\subsection{Roll-up under an end moment}
\label{sec:rollup}

A cantilever under $M = 2\pi EI/L$ must roll into an exact closed circle, so the
reference requires no lookup at all.

\begin{figure}[H]
\centering
\includegraphics[width=\textwidth]{rollup.png}
\caption{Roll-up at 30, 60 and 100\% of $2\pi EI/L$, and the distance from the
exact circular arc. Billow converges at second order; \code{kite\_fem} moves
away from the exact answer under refinement.}
\label{fig:rollup}
\end{figure}

\begin{table}[H]
\centering
\caption{Tip distance from the exact arc, normalised by length.}
\begin{tabular}{lrrr}
\toprule
Elements & 10 & 20 & 40 \\
\midrule
Billow (4 load steps) & $1.13\times10^{-1}$ & $3.16\times10^{-2}$ & $8.14\times10^{-3}$\\
\code{kite\_fem} (single shot) & $5.90\times10^{-2}$ & $1.71\times10^{-1}$ & $2.32\times10^{-1}$\\
\bottomrule
\end{tabular}
\end{table}

Billow's error ratios are $3.6$ and $3.9$, i.e.\ second order. Extending to 80
elements gives $2.05\times10^{-3}$, confirming the trend. This benchmark is what
exposed the moment-conjugacy defect of \S\ref{sec:moments}.

\subsection{Bathe--Bolourchi \texorpdfstring{$45^\circ$}{45-degree} bend}

The standard three-dimensional finite-rotation benchmark: a $45^\circ$ circular
arc of radius 100, square section, out-of-plane tip load $F = 600$. Published
tip \emph{displacements} differ by about $1\%$ between authors, so the reference
is a band.

\begin{figure}[H]
\centering
\includegraphics[width=\textwidth]{bend45.png}
\caption{Tip displacement against published references, and convergence.}
\label{fig:bend45}
\end{figure}

\begin{table}[H]
\centering
\caption{Tip displacement, 32 elements. The arc here starts at the origin with
its tangent along $+x$ and sweeps towards $+y$, which swaps the two in-plane
components relative to the usual tabulation.}
\begin{tabular}{lrrrr}
\toprule
& $u_x$ & $u_y$ & $u_z$ & $\lVert\mathbf{u}\rVert$\\
\midrule
Simo \& Vu-Quoc (1986) & $-23.48$ & $-13.50$ & $53.37$ & $59.85$\\
Bathe \& Bolourchi (1979) & $-23.5$ & $-13.4$ & $53.4$ & $59.85$\\
Crisfield (1990) & $-23.5$ & $-13.6$ & $53.4$ & $59.88$\\
\midrule
\textbf{Billow} & $\mathbf{-23.52}$ & $\mathbf{-13.60}$ & $\mathbf{53.44}$ & $\mathbf{59.95}$\\
\code{kite\_fem} & $-37.74$ & $-22.14$ & $28.17$ & $52.04$\\
\bottomrule
\end{tabular}
\end{table}

Billow sits within the spread between published references ($0.17\%$ on the
magnitude). \code{kite\_fem} is $47\%$ low out of plane.

\subsection{Inflatable tube law}

The port of \S\ref{sec:inflatable} is verified with \textbf{pure end moments},
where the exact solution is a constant-curvature arc at whatever curvature the
fit prescribes. That check is exact at any deflection --- unlike the tip-load
form the fit was published in, whose own inversion assumes linear cantilever
theory and therefore cannot serve as a verification instrument.

\begin{figure}[H]
\centering
\includegraphics[width=\textwidth]{inflatable.png}
\caption{ASKITE fits (lines) with Billow end-moment solves on top (markers).
Hollow markers lie beyond the collapse curvature, where the fit is
extrapolation: the solve is fine, the calibration is not.
This figure completes \code{kite\_fem/examples/FEM\_beam\_verification.py},
which plots the fitted curves but leaves the simulation loop as comments.}
\label{fig:inflatable}
\end{figure}

Worst errors: $0.031\%$ (bending, 0.3\,bar), $0.057\%$ (bending, 0.5\,bar),
$0.749\%$ and $1.491\%$ (torsion).

\subsection{Beam convergence and shear locking}
\label{sec:locking}

Tip deflection of a slender cantilever must approach
$PL^3/3EI + PL/\kappa GA$ at second order:

\begin{table}[H]
\centering
\begin{tabular}{lrrrrrr}
\toprule
Elements & 2 & 4 & 10 & 20 & 40 & 80\\
\midrule
Error vs Timoshenko & $-6.25\%$ & $-1.56\%$ & $-0.250\%$ & $-0.063\%$
 & $-0.016\%$ & $-0.004\%$\\
\bottomrule
\end{tabular}
\end{table}

The error falls by a factor of four per refinement. Locking is checked
separately by holding the load and lengthening the beam through slenderness
ratios of 50, 200 and 800: the tip deflection tracks beam theory to $0.2\%$
throughout, where a locking element would stiffen towards zero.

\subsection{Membrane}

Checked against closed form rather than against another code:

\begin{itemize}
\item \textbf{Rigid-body motion} costs zero energy under arbitrary rotation and
translation --- what the finite-strain measure buys over a small-strain one.
\item \textbf{Equibiaxial stretch} matches the analytic SVK plane-stress energy
to $10^{-9}$ relative.
\item \textbf{Wrinkled branch} is exactly uniaxial: stretching along $x$ while
squashing along $y$ gives $\tfrac{1}{2}E t A e_1^2$, and squashing further
changes nothing.
\item \textbf{Slack panel} stores exactly zero, where the unrelaxed law would
store large energy.
\item \textbf{Flat-plate inflation} converges under mesh refinement from a
configuration with zero out-of-plane stiffness --- the singular start that a
residual-form Newton method has nothing to invert at.
\end{itemize}

\subsection{Test suite}

86 tests, all passing, organised one file per concern:

\begin{table}[H]
\centering
\begin{tabular}{ll}
\toprule
\code{test\_rotations.py} & Cayley round trips, half-angle identity\\
\code{test\_cable.py} & Hooke parity with PSS, slack cut, pulley tension equalisation\\
\code{test\_beam.py} & objectivity, $O(h^2)$ convergence, no locking, torsion vs $GJ$\\
\code{test\_membrane.py} & analytic SVK, three wrinkling branches, inflation, scale\\
\code{test\_inflatable.py} & energy integrates the fitted moment; end-moment solves\\
\code{test\_benchmarks.py} & roll-up vs exact circle, load-step independence, bend45\\
\code{test\_energy.py} & mapped assembly vs a naive per-element sum; DOF layout\\
\bottomrule
\end{tabular}
\end{table}

The most important of these is in \code{test\_energy.py}: the total energy
computed through the gather-and-map machinery is compared against the same
energy assembled one element at a time in plain Python, on a deliberately
irregular model containing every element type with a beam running over a
scattered, out-of-order subset of the node numbering. That machinery is easy to
get subtly wrong --- a column-major reshape, a rotational slot off by one ---
and wrong in a way that still converges to a plausible-looking shape.

%======================================================================
\section{Demonstration cases}

\subsection{Wrinkling: fabric in a frame smaller than itself}

\begin{figure}[H]
\centering
\includegraphics[width=\textwidth]{wrinkling.png}
\caption{A panel clamped into a frame 6\% smaller than itself, then pushed out
of plane. The relaxed law converges in 31 iterations to a smooth billow; the
unrelaxed one takes 211 and buckles into mesh-scale crumple at less than
two-thirds the deflection.}
\label{fig:wrinkling}
\end{figure}

This is the clearest illustration of why the relaxed energy is not optional.
The unrelaxed panel stores $40.6$\,J against the relaxed panel's $0.116$\,J,
almost all of it fighting compression that fabric cannot carry.

\subsection{Large-deflection cantilever}

\begin{figure}[H]
\centering
\includegraphics[width=\textwidth]{cantilever.png}
\caption{Deflected shapes under a tip load swept from 50 to 900\,N, and tip
deflection against linear theory. The beam visibly shortens in $x$ as it bends
--- the geometric nonlinearity linear theory omits. At the largest load, linear
theory overshoots by 51\%.}
\label{fig:cantilever}
\end{figure}

Median 9 IPOPT iterations per load step with warm starting.

\subsection{Batten-stiffened sail}

\begin{figure}[H]
\centering
\includegraphics[width=\textwidth]{sail.png}
\caption{Beams, membranes, cables and a trailing-edge hem in one energy: a
spar carrying bending and torsion, a wrinkling canopy hanging off it, three
bridles to a fixed anchor. 393 DOF, 54 iterations, 1.4\,s, residual
$1.0\times10^{-9}$\,N.}
\label{fig:sail}
\end{figure}

Note that only 13 of the nodes carry a material frame; the canopy nodes stay at
three degrees of freedom.

\subsection{Scaling}

\begin{figure}[H]
\centering
\includegraphics[width=\textwidth]{scaling.png}
\caption{Build and solve time against problem size, and iteration count.
Because the graph is built once per element \emph{type}, build cost tracks the
solve rather than the element count, and the iteration count is nearly mesh
independent across a 37$\times$ range in degrees of freedom.}
\label{fig:scaling}
\end{figure}

%======================================================================
\section{Comparison with \texorpdfstring{\code{kite\_fem}}{kite\_fem}}

\code{kite\_fem} (with \code{pyfe3d}) is the existing FEM path in the AWETrim
tree. This section states what differs, what was measured, and where the
comparison is and is not fair.

\subsection{Capability}

\begin{table}[H]
\centering
\small
\begin{tabular}{p{4.0cm}p{5.1cm}p{5.1cm}}
\toprule
& \code{kite\_fem} / \code{pyfe3d} & \textbf{Billow}\\
\midrule
Elements & springs, pulleys, inflatable Timoshenko beams
         & cables, pulleys, geometrically exact Timoshenko beams,
           inflatable tubes, wrinkling membranes\\
Canopy fabric & noncompressive spring net (no membrane element exists)
              & CST membrane, Pipkin relaxed energy\\
Degrees of freedom & 6 per node everywhere, masked via \code{bc}
                   & 3 per node, $+3$ only on beam nodes\\
Formulation & Newton--Raphson on $\mathbf{f}_e-\mathbf{f}_i$
            & minimum potential energy (IPOPT)\\
Assembly & scipy sparse COO per element & one mapped kernel per element type\\
Regularisation & \code{I\_stiffness=25}, \code{step\_limit=0.2},
                 \code{relax=0.5}, optional pseudo-transient
               & IPOPT inertia correction; membrane slack blend\\
Warm start & none (\code{solve} resets to undeformed) & state and duals\\
Default tolerance & $10^{-2}$\,N & $10^{-8}$ IPOPT, accepted at $10^{-6}$\,N\\
Inflatable tube & secant $EI$/$GJ$ plus collapse flag
                & the same fits integrated to $W(\kappa)$\\
Experimental validation & hanging-kite test against measured data & none\\
\bottomrule
\end{tabular}
\end{table}

\subsection{Two traps in making it a fair comparison}

\begin{enumerate}
\item \code{FEM\_structure.solve} defaults to \code{I\_stiffness=25}, an
\textbf{absolute} regularisation in N/m added to every degree of freedom of the
tangent. The benchmark beams here have $EI/L^3 = 10$\,N/m, so the default is
larger than the structure: 2000 iterations with the residual stuck at 20\,N.
With it at zero the same problem converges in 38. It is a kite-scale tuning
constant, not a universal one. All results above use zero.
\item The property matching was verified in the linear limit before any
conclusion was drawn: below about one degree of rotation the two codes agree to
$0.02\%$.
\end{enumerate}

\subsection{Cost and accuracy on a common problem}

Tip-loaded slender cantilever, matched properties, same mesh, one machine, one
run.

\begin{table}[H]
\centering
\small
\caption{Source: \code{run\_cost\_comparison.py}. ``Warm'' is a second solve at
a 5\% different load; \code{kite\_fem} restarts from the undeformed
configuration on every call, so its warm column is a second cold solve.}
\begin{tabular}{rrlrrrrrr}
\toprule
$\alpha$ & El. & Model & DOF & Setup (s) & Cold (s) & Warm (s) & Iter. & ms/iter\\
\midrule
1 & 20 & Billow           & 126 & 0.133 & 0.249 & \textbf{0.081} & 16 & 15.6\\
1 & 20 & \code{kite\_fem} & 126 & \textbf{0.004} & \textbf{0.132} & 0.151 & 39 & \textbf{3.4}\\
1 & 80 & Billow           & 486 & 0.217 & 0.733 & \textbf{0.282} & 17 & 43\\
1 & 80 & \code{kite\_fem} & 486 & \textbf{0.009} & \textbf{0.189} & 0.192 & 39 & \textbf{4.8}\\
6 & 40 & Billow           & 246 & 0.076 & \textbf{0.333} & \textbf{0.062} & \textbf{22} & 15\\
6 & 40 & \code{kite\_fem} & 246 & \textbf{0.005} & 1.384 & 1.526 & 444 & \textbf{3.1}\\
6 & 80 & Billow           & 486 & 0.186 & \textbf{0.669} & \textbf{0.233} & \textbf{23} & 29\\
6 & 80 & \code{kite\_fem} & 486 & \textbf{0.005} & 2.270 & 2.658 & 490 & \textbf{4.8}\\
\bottomrule
\end{tabular}
\end{table}

\code{kite\_fem} has roughly $6\times$ cheaper iterations (compiled \code{pyfe3d}
plus scipy sparse) and $30\times$ faster setup. Billow uses 2--20$\times$ fewer
iterations, and its count is nearly independent of both mesh and load
(16--23 throughout) where \code{kite\_fem}'s climbs from 39 to about 490 as the
load grows. \code{kite\_fem} therefore wins on easy, small, one-shot problems;
Billow wins as soon as the problem is hard, reused, or large. Billow's setup
cost is a one-off amortised across a sweep --- exactly the coupled-loop pattern.

\subsection{The canopy model: spring net versus membrane}

Since \code{kite\_fem} has no membrane element, its canopy is a noncompressive
spring net (chordwise and spanwise links, \emph{no diagonals}; its own example
carries a commented-out block of ``fictional springs constraining the canopy from
collapse''). Billow can build exactly that net from its cable elements, which
makes a genuine like-for-like study possible.

\begin{figure}[H]
\centering
\includegraphics[width=\textwidth]{canopy_model.png}
\caption{Spring net versus CST membrane on identical meshes, both minimised.
Left: both converge under out-of-plane load, to answers 7\% apart. Middle: under
equibiaxial stretch the net is too soft by exactly $1-\nu$. Right: under in-plane
shear the membrane energy scales as $\gamma^{2.0}$ and the net as $\gamma^{4.0}$
--- the net has \emph{no shear stiffness at leading order}.}
\label{fig:canopymodel}
\end{figure}

Measured: biaxial energy ratio $0.728$ against $1-\nu = 0.70$; shear exponents
$2.018$ (membrane) and $3.999$ (net).

\subsubsection{Taut: shape and cost}

\begin{figure}[H]
\centering
\includegraphics[width=\textwidth]{taut_net_vs_membrane.png}
\caption{Taut bay, $24\times18$, both minimised on the same mesh with the same
1425 degrees of freedom, so this is purely the canopy model. The net billows
$9.6\%$ deeper; the best single scale factor is $1.120$, and a $3.31$\,mm
residual shape difference remains after removing it.}
\label{fig:tautshape}
\end{figure}

The local softness ratio is not constant: $1.10$ at mid-chord rising to $1.145$
near the clamped edges, where the fabric carries biaxial tension and shear that
the net cannot represent. The residual map shows the net relatively too deep in
the corners and too shallow in the middle --- it billows into a slightly
different surface, not merely a deeper one.

\textbf{The scale factor cannot be calibrated away}, because it depends on the
stress state: about $1.10$--$1.12$ when taut, but $1.00$ when fully wrinkled
(the slack cases agree within $1\%$, since a deeply slack panel carries
essentially uniaxial chordwise tension where the two models coincide). The
correction needed would change with flight condition.

Cost is not a differentiator. At $24\times18$: net $0.092$\,s setup and
$0.227$\,s solve in 12 iterations; membrane $0.176$\,s and $0.245$\,s in 8. The
membrane's per-iteration cost is higher but it needs fewer iterations, and the
two roughly cancel.

\subsubsection{Robustness}

\begin{table}[H]
\centering
\small
\caption{Peak billow (\% chord) and section RMS against each model's own finest
mesh. The net fails outright at one intermediate slack mesh --- the signature of
its zero-energy shear mode. The membrane never failed.}
\begin{tabular}{lrrrrrr}
\toprule
& \multicolumn{3}{c}{Spring net} & \multicolumn{3}{c}{Membrane}\\
\cmidrule(lr){2-4}\cmidrule(lr){5-7}
Mesh & Camber & RMS & Conv. & Camber & RMS & Conv.\\
\midrule
\multicolumn{7}{l}{\emph{taut}}\\
$4\times3$   & 2.858 & 0.388 & yes & 2.542 & 0.401 & yes\\
$8\times6$   & 3.129 & 0.051 & yes & 2.842 & 0.051 & yes\\
$16\times12$ & 3.117 & 0.010 & yes & 2.842 & 0.010 & yes\\
$24\times18$ & 3.115 & ---   & yes & 2.842 & ---   & yes\\
\multicolumn{7}{l}{\emph{slack 3\%}}\\
$8\times6$   & 7.388 & 0.186 & yes & 7.449 & 0.143 & yes\\
$12\times9$  & 7.240 & 0.091 & yes & 7.294 & 0.097 & yes\\
$16\times12$ & 5.675 & 1.719 & \textbf{no} & 7.412 & 0.031 & yes\\
$24\times18$ & 7.302 & ---   & yes & 7.402 & ---   & yes\\
\bottomrule
\end{tabular}
\end{table}

\subsubsection{Canopy model versus solver, separated}

A single bay isolates the two effects that otherwise get conflated. Rows 2 and 3
solve the \emph{identical} spring net, so their difference is the solver alone.

\begin{table}[H]
\centering
\small
\caption{One $2.0\times1.3$\,m bay at 250\,Pa, all edges fixed.}
\begin{tabular}{llrrrr}
\toprule
Case & Configuration & Solve (s) & Iterations & Camber (\% c) & Converged\\
\midrule
taut $8\times6$ & Billow + membrane   & 0.030 & 7 & 2.842 & yes\\
taut $8\times6$ & Billow + net        & 0.026 & 10 & 3.129 & yes\\
taut $8\times6$ & \code{kite\_fem} + net & 10.04 & 2729 & 3.129 & yes\\
\midrule
taut $12\times9$ & Billow + net        & 0.034 & 9 & 3.083 & yes\\
taut $12\times9$ & \code{kite\_fem} + net & 22.00 & 3001 & 5.283 & \textbf{no}\\
\midrule
slack $8\times6$ & Billow + net        & 0.049 & 23 & 7.388 & yes\\
slack $8\times6$ & \code{kite\_fem} + net & 9.06 & 3001 & 7.363 & \textbf{no}\\
\bottomrule
\end{tabular}
\end{table}

Where both converge they agree exactly ($3.129\%$ both), which is what makes the
\code{kite\_fem} driver used throughout this document trustworthy. On the
identical net, minimum energy is roughly $390\times$ faster than the damped
Newton solve, and remains convergent where the Newton solve does not.

\subsection{Full kite}

A V3-scale structure --- inflated leading-edge tube, inflated struts at every bay
boundary, canopy over nine bays, trailing-edge hem, bridle to a fixed anchor,
uniform 250\,Pa --- was built with the same nodes, tubes and cables in both codes.

\begin{table}[H]
\centering
\small
\begin{tabular}{lrrrrrl}
\toprule
Model & DOF & Setup (s) & Cold (s) & Warm (s) & Iter. & Residual (N)\\
\midrule
\textbf{Billow} & 1893 & 0.61 & \textbf{5.91} & \textbf{0.58} & 73
 & $7.2\times10^{-9}$ (converged)\\
\code{kite\_fem} ($I{=}25$) & 2976 & 0.05 & 76.2 & 72.0 & 2001 (cap)
 & $4.1\times10^{5}$ (diverged)\\
\code{kite\_fem} ($I{=}0$) & 2976 & 0.06 & 108.9 & 68.9 & 2001 (cap)
 & $1.3\times10^{8}$ (diverged)\\
\bottomrule
\end{tabular}
\end{table}

\paragraph{Interpretation, carefully.} No usable \code{kite\_fem} full-kite
number was obtained. Diagonals, canopy curvature, both regularisation settings
and load stepping with rebuilt initial conditions were all tried; all diverged,
with residuals roughly $100\times$ the applied load. \textbf{This is consistent
with a known open issue on the \code{kite\_fem} side} rather than a defect in the
driver used here --- the ASKITE integration commit records ``kite\_fem still
diverging, but not due to the coupling''. The only kite-scale figure that comes
from its own workflow is its hanging-kite validation, at 442--671\,s per load
case. The single-bay comparison above is the trustworthy cost comparison.

%======================================================================
\section{How many membrane triangles?}

Chordwise and spanwise resolution are swept independently, because they are not
the same problem: a canopy bay billows across the chord --- which sets the
section the aerodynamics sees --- while spanwise it varies slowly between struts.

The test case is one bay of the LEI-V3 canopy: chord 2.0\,m, strut spacing
1.3\,m (the V3 spans 8.28\,m over 10 stations), held on all four edges, under
250\,Pa. Both a taut panel and a realistically slack one (reference cut 3\%
oversize) are run.

Convergence is judged on quantities the coupling actually consumes: peak billow
over chord, and the RMS deviation of the mid-span chordwise profile from the
finest mesh. \textbf{Peak local strain is deliberately not a criterion}: in a
wrinkling membrane it does not converge, and nothing downstream reads it.

\begin{figure}[H]
\centering
\includegraphics[width=\textwidth]{mesh_shapes.png}
\caption{One bay under uniform refinement. A $2\times2$ bay is a tent, not a
canopy: the camber is only 2.8\% low but the section is two straight lines.
By $8\times8$ the profile is visually converged.}
\label{fig:meshshapes}
\end{figure}

\begin{figure}[H]
\centering
\includegraphics[width=\textwidth]{mesh_requirement.png}
\caption{Camber and section-shape convergence, chordwise (solid) and spanwise
(dashed), and cost per bay.}
\label{fig:meshreq}
\end{figure}

\begin{table}[H]
\centering
\caption{Candidate meshes against a $24\times24$ reference, slack 3\%.}
\begin{tabular}{lrrrrr}
\toprule
Mesh & Triangles & Camber error & Section RMS (\% c) & Build $+$ solve (s) & Iter.\\
\midrule
$4\times4$   & 32  & 0.88\% & 0.967 & 0.03 & 12\\
$6\times6$   & 72  & 0.61\% & 0.338 & 0.07 & 17\\
\textbf{$8\times6$} & \textbf{96} & \textbf{0.71\%} & \textbf{0.145} & \textbf{0.09} & 19\\
$8\times8$   & 128 & 0.16\% & 0.140 & 0.17 & 15\\
$12\times8$  & 192 & 0.18\% & 0.075 & 0.18 & 21\\
$16\times12$ & 384 & 0.21\% & 0.033 & 0.44 & 22\\
\bottomrule
\end{tabular}
\end{table}

\paragraph{Recommendation.} $\mathbf{8}$ \textbf{chordwise} $\times$
$\mathbf{6}$ \textbf{spanwise per bay}, 96 triangles, about 900 for the whole
nine-bay canopy. Camber error $0.71\%$ and section RMS $0.145\%$ of chord ---
2.9\,mm on a 2\,m chord, comfortably below what an airfoil parameterisation
resolves. Use $12\times8$ for final runs: it halves the shape error for twice
the cost.

\paragraph{Slack costs roughly twice the mesh.} At 3\% oversize the panel is
88\% wrinkled (against 1--4\% taut) and needs about $12\times8$ for the same
accuracy, with 4$\times$ the iterations. Since a real canopy flies slack, the
as-cut slack of the actual sail drives the mesh requirement more than anything
else here and is worth pinning down.

\paragraph{Caveat.} These numbers are for \emph{uniform} pressure. A real
chordwise pressure distribution is peaked near the leading edge and may demand
more chordwise resolution; the study should be repeated once loads come from an
aerodynamic solve.

%======================================================================
\section{Known limitations}

\begin{description}[style=unboxed,leftmargin=1.2em]
\item[No experimental validation.] Every check in this document is against
closed-form solutions, published benchmarks, or the ASKITE fits. The claim that
the membrane is ``more correct'' than a spring net rests on it being the
physically appropriate model, not on measurement. \code{kite\_fem} has a
hanging-kite test against real data; Billow has nothing equivalent yet.

\item[Applied moments.] Need load stepping, and are exact only about a fixed
axis; a dead three-dimensional moment is non-conservative and has no potential.
Nodal forces are unaffected.

\item[Post-collapse tubes.] Reported, not modelled. The fits cover bending and
torsion only; axial and shear stiffness must be supplied from tube geometry.

\item[Dead loads.] Loads are frozen per solve, which matches a staggered
coupling loop. A true follower pressure would not be conservative and could not
be posed as a potential --- worth remembering before anyone tries to fold
aerodynamics into the same minimisation.

\item[No contact.] Contact and self-contact are not modelled.

\item[Setup cost.] About $0.2$\,s of graph construction for a 486-DOF beam is
disproportionate against \code{kite\_fem}'s $0.005$\,s, even though it is a
one-off and stops mattering at membrane scale.

\item[Threaded assembly.] \code{PotentialEnergy(parallelization=...)} exposes
CasADi threaded maps but is untested at scale; serial is the default.

\item[Not coupled.] There is no aerodynamic coupling adapter yet.
\end{description}

%======================================================================
\section{Reproducing everything in this document}

\begin{table}[H]
\centering
\small
\begin{tabular}{ll}
\toprule
Figure / table & Command\\
\midrule
Figs.~\ref{fig:wrinkling}--\ref{fig:scaling}, \ref{fig:canopymodel}
 & \code{python scripts/structural/run\_demo\_cases.py}\\
Figs.~\ref{fig:elastica}--\ref{fig:inflatable}
 & \code{python scripts/structural/run\_validation\_benchmarks.py}\\
Cantilever cost table & \code{python scripts/structural/run\_cost\_comparison.py}\\
Figs.~\ref{fig:meshshapes}, \ref{fig:meshreq}
 & \code{python scripts/structural/run\_mesh\_requirement.py}\\
Single-bay comparison & \code{python scripts/structural/run\_panel\_cost.py}\\
Full-kite table & \code{python scripts/structural/run\_full\_kite\_cost.py}\\
Test suite & \code{pytest tests/structural/}\\
\bottomrule
\end{tabular}
\end{table}

\subsection{Package layout}

\begin{table}[H]
\centering
\small
\begin{tabular}{ll}
\toprule
\code{rotations.py} & $SO(3)$ kernels over an \code{xp} namespace (NumPy or CasADi)\\
\code{model.py} & \code{DofLayout}, \code{StructuralState}, \code{StructuralModel}\\
\code{energy.py} & \code{PotentialEnergy}: mapped assembly, gradient, tangent stiffness\\
\code{solver.py} & \code{MinimumEnergySolver}, \code{StructuralSolution}\\
\code{elements/base.py} & \code{ElementKernel} protocol, \code{ElementSet}\\
\code{elements/cable.py} & \code{CableKernel}, \code{PulleyKernel}\\
\code{elements/beam.py} & \code{TimoshenkoBeamKernel}, \code{BeamSection}\\
\code{elements/inflatable.py} & \code{InflatableTubeLaw}, \code{InflatableBeamKernel}\\
\code{elements/membrane.py} & \code{MembraneKernel}, \code{membrane\_regimes}\\
\bottomrule
\end{tabular}
\end{table}

\subsection{Minimal example}

\begin{quote}\small\ttfamily
\noindent
import numpy as np\\
from awetrim.structural import MinimumEnergySolver, StructuralModel\\
from awetrim.structural.elements import build\_cable\_elements\\[4pt]
nodes = np.array([[0., 0, 0], [1., 0, 0]])\\
cables = build\_cable\_elements([[0, 1]], rest\_lengths=[1.0], stiffness=[1e4])\\
model = StructuralModel(nodes, [cables], fixed\_translation\_nodes=[0])\\
solution = MinimumEnergySolver(model).solve(np.array([[0., 0, 0], [10., 0, 0]]))
\end{quote}

\subsection{Invariants to preserve}

\begin{itemize}
\item \textbf{Isolation.} NumPy and CasADi only.
\item \textbf{One mapped kernel per element type.} Never reintroduce a Python
loop over elements that builds \code{MX}.
\item \textbf{Single source for every formula.} \code{rotations.py} holds the
$SO(3)$ maps against an \code{xp} namespace so NumPy and CasADi share one
implementation; \code{beam\_strains} and \code{green\_strain} are called by both
the kernels and the NumPy setup and diagnostic paths.
\item \textbf{No CasADi across the module boundary.} Symbolics are created inside
\code{energy.py} and stay there.
\item \textbf{Kernels are pure.}
\item \textbf{\code{converged} is the physical verdict}, with IPOPT's own kept
alongside.
\end{itemize}

%======================================================================
\section*{References}
\addcontentsline{toc}{section}{References}

\begin{enumerate}[label={[\arabic*]}]
\item J.C. Simo and L. Vu-Quoc, ``A three-dimensional finite-strain rod model.
Part II: computational aspects'', \emph{Comput. Methods Appl. Mech. Engrg.}
\textbf{58} (1986) 79--116.
\item A. Ibrahimbegovi\'c, ``On the choice of finite rotation parameters'',
\emph{Comput. Methods Appl. Mech. Engrg.} \textbf{149} (1997) 49--71.
\item K.-J. Bathe and S. Bolourchi, ``Large displacement analysis of
three-dimensional beam structures'', \emph{Int. J. Numer. Meth. Engng.}
\textbf{14} (1979) 961--986.
\item A.C. Pipkin, ``The relaxed energy density for isotropic elastic
membranes'', \emph{IMA J. Appl. Math.} \textbf{36} (1986) 85--99.
\item D.G. Roddeman, J. Drukker, C.W.J. Oomens and J.D. Janssen, ``The wrinkling
of thin membranes'', \emph{J. Appl. Mech.} \textbf{54} (1987) 884--892.
\item K.E. Bisshopp and D.C. Drucker, ``Large deflection of cantilever beams'',
\emph{Q. Appl. Math.} \textbf{3} (1945) 272--275.
\item Y. Luo, ``An efficient 3D Timoshenko beam element with consistent shape
functions'', \emph{Adv. Theor. Appl. Mech.} \textbf{1}(3) (2008) 95--106.
\item O. Cayon, M. Gaunaa and R. Schmehl, ``Fast aero-structural model of a
leading-edge inflatable kite'', \emph{Energies} \textbf{16} (2023) 3061.
\item ASKITE, \url{https://github.com/awegroup/ASKITE};
\code{kite\_fem}, \url{https://github.com/awegroup/kite_fem}.
\end{enumerate}

\end{document}
